% ===== §6 The Ethics of Reticence =====
\section{The Ethics of Reticence}
\label{p18:sec:ethics}

\noindent
The criterion of \xref{p18:sec:criterion} establishes that effacement is a harm; the construction of \xref{p18:sec:reticence} establishes what reticence is and how it repairs the harm. Neither yet establishes that reticence is \emph{owed}---that non-determination is a requirement of justice and not a mere kindness, an aesthetic preference for relations in which both flourish. This section supplies the ground. It moves first negatively, showing that each major ethical framework registers something of the wrong yet, pressed to ground the obligation, fails at a diagnostic point of its own; then positively, locating the ground in the standing of the other as a generative being; and finally reflexively, establishing the constraint that forbids the ethic its own paternalistic application. The negative movement is not a survey but a series of defeats, and the most instructive defeat is care ethics', because care ethics comes nearest and fails most revealingly.

\subsection{Four frameworks, four failures}
\label{p18:sub:eth-insufficiency}

\emph{Consequentialism cannot see the harm.} An ethic evaluating by outcomes is structurally blind here, for the outcomes of effacement are in aggregate good: value rises, the relation prospers, both may report contentment. A framework that sums welfare cannot find a harm invisible in the sum, and effacement is such a harm, since what it destroys---the weaker pole's own generativity---is replaced in the ledger by the stronger pole's surplus, leaving the total undiminished. This is not a marginal failure of detail but the sharpest demonstration that the wrong is non-aggregative.

\emph{Contractualism is defeated by consent.} An ethic locating wrong in what could not be justified to, or was not agreed by, the other founders on the datum of \xref{p18:sub:phen-sincere}: the effaced party consents, endorses, is grateful, would choose it again. Contractualism can condemn a determination imposed over objection; it has nothing to say of one the other welcomes. Yet the harm is present in the welcoming case, because consent operates at the level of the will and effacement at the level of being. A wrong surviving the genuine consent of the one wronged lies outside the reach of any ethic that makes consent its criterion.

\emph{Kantian respect locates the wrong too high.} The deontological tradition forbids treating another merely as a means and requires treating her always also as an end---and fixes the end-status in her \emph{rational agency}, her capacity to set and pursue ends. But rational agency survives effacement intact: the effaced party still reasons, still chooses, still wills, and often wills her own effacement. An ethic that secures her standing as a rational chooser secures exactly what effacement leaves untouched. The Kantian framework is not wrong; it aims at the right relation (treating-as-end) but at the wrong faculty, and so passes over the harm. (The correction---relocating the end-status from rational to generative agency---is the work of \xref{p18:sub:eth-ground} below.)

\emph{Care ethics supplies the ontology and valorises the harm.} Of the four, care ethics alone begins from the relational ontology this paper requires: it knows selves are constituted in and sustained by relation, and attends to dependency and to the labour of maintenance the other frameworks overlook. Noddings builds the ethical relation from the ``one-caring'' and the ``cared-for,'' and makes its central act \emph{engrossment}---the one-caring's apprehension of the other's reality, the displacement of her own motives to act on the other's behalf \parencite{noddings1984}. Here, if anywhere, an ethic should catch a harm done through devotion. Instead care ethics, in its received form, \emph{valorises} the very structure that effaces. For engrossment and motivational displacement, raised to the paradigm of the ethical, describe with precision the inner life of the effacing pole: the one who displaces her own generativity into the other's, who is engrossed in the other to the point of generating only for and as the other. The standing feminist critique of care ethics already saw the danger from the side of the carer---that the ethic celebrates self-sacrifice and risks the caregiver's exploitation, pressing her toward a caring that forgets her own needs. This paper locates the symmetrical danger on the other side: that the \emph{cared-for} can be effaced by an engrossment too complete, generated-for so thoroughly that she ceases to generate. Care ethics supplies the indispensable relational ontology, and must be corrected at exactly the point where it makes a virtue of the unstinting giving that, past a threshold the criterion of \xref{p18:sec:criterion} makes exact, occupies the other out of her own being.

\begin{claim}[Why the frameworks fail]
Each major framework fails to ground the obligation of reticence at a diagnostic point: consequentialism because the harm is non-aggregative and invisible in the sum; contractualism because the harm survives sincere consent; Kantian respect because it secures the rational agency effacement leaves intact; care ethics because it valorises, as devotion, the engrossment that effaces. The pattern of failure is itself informative. Three frameworks aim too high---at outcomes, at the will, at rational agency---missing a harm that lies beneath outcome, will, and reason, at the level of generative being; and the one framework pitched low enough, care ethics, mistakes the harm for the good. The ground of reticence must therefore be sought at the level the first three overshoot and the fourth misreads.
\end{claim}

\subsection{The positive ground: the standing of the generative being}
\label{p18:sub:eth-ground}

If reticence is owed, it is owed because the other's standing as a generative being is a good making a direct claim---a claim damaged specifically by the determination of her generative domain, and damaged even when the determination is benevolent and welcomed. The ground has three components, each drawn from a tradition and each corrected where its received form would miss the harm.

\emph{The corrected end-formula.} Relocate the Kantian end-status from rational to \emph{generative} agency: the other's standing as a being who continually constitutes herself by generating at her own frequency. The formula becomes---act so that you never determine the other's generative domain merely as an extension of your own. To overwrite her frequency with yours, however lovingly, is to treat her generativity as material for your own, making of her generative being a means to the relation's flourishing rather than an end whose exercise the relation exists to sustain. This repairs the defect diagnosed in \xref{p18:sub:eth-insufficiency}: the end-formula was aimed at the wrong faculty, and aimed right, it reaches the harm.

\emph{The foreclosure of activity.} From the tradition that locates a being's good in the \emph{exercise} of its characteristic activity---its \emph{energeia}---rather than in any state or resource, the harm comes into focus as the foreclosure of an activity rather than the withholding of a good. What the effaced party loses is not something she has but something she \emph{does}: the activity of generating, in which her flourishing consists. A capability account at the level of functioning catches part of this, but the harm lies one level deeper than a lost functioning---it is the loss of the self-constituting activity any particular functioning would express. Reticence, in this register, is the virtue that holds open the field of another's exercise.

\emph{The second-personal warrant.} To the question \emph{on whose authority} the claim is made, the answer is that the obligation is second-personal---owed to her, from her standpoint, on her own standing to demand it \parencite{darwall2006}. The warrant is not a third-personal fact about the goodness of generativity assessed from outside, but her own, perhaps unexercised, standing to require that her generative domain not be determined for her. This component prepares the constraint of \xref{p18:sub:eth-paternalism}: if the warrant is hers, the remedy cannot be administered over her head without reproducing the wrong.

\subsection{A duty grounded, a virtue in application}
\label{p18:sub:eth-virtue}

The ground is deontological: reticence is a negative obligation owed in virtue of the other's standing. But an obligation of this kind cannot be discharged by rule, and the reason lies in the harm's structure. Effacement arrives through devotion, in the ordinary motions of a loving and capable subject doing well by the relation; it breaks no rule and feels, from inside, like care. One cannot rule-follow one's way clear of it, because no rule fires---there is no moment the stronger pole can recognise, by rule, as the wrong. What is required instead is a trained \emph{perception}: the cultivated capacity to notice, amid generous and well-meant generation, that one's own activity has begun to foreclose another's. The ethic is thus deontological in its ground and virtue-theoretic in its application, and the two are compatible: the obligation is fixed and negative, its \emph{detection} a matter of cultivated sensibility no algorithm supplies. This is not a retreat into vagueness but a recognition of the harm's phenomenology---a wrong that announces itself by no transgression can be caught only by a perception trained to its signs.

\subsection{The reflexive limit: reticence without paternalism}
\label{p18:sub:eth-paternalism}

The ethic carries a danger of its own making, and must close by disarming it. A stronger pole, having grasped the analysis, may set out to \emph{judge} that the other is effaced and to \emph{act} to restore her generativity---to manage her back into generation. But to judge another's effacement on her behalf and take her restoration in hand is itself a determination of her domain: it sets her a project, her own recovery, she did not constitute, effacing her in a second key, now in the name of repair. The danger is not hypothetical; it is the precise form that care ethics' critics already identified within devotion itself. Goodin's worry about care---that subsuming individuals into a relationship of ``we''-ness risks losing track of the separateness of persons, issuing in a smothering paternalism \parencite{goodin1985}---names exactly the recursion that threatens here: an ethic of non-determination, applied as management, smothers in the act of rescuing.

The constraint that disarms it follows from the second-personal ground. Reticence is the withdrawal of one's \emph{own} determining activity, never the management of the other's. The stronger pole's obligation extends to vacating the domain it has occupied, and stops there; it does not extend to administering what the other then does with the vacated room, for that administration re-occupies the room in the act of giving it. Where the very question is whether the other is effaced, the only non-effacing route is to return the question to her---to leave her reading of her own generativity undetermined---which is itself an instance of the first channel of \xref{p18:sub:ret-channels}, relinquished now at the meta-level.

\begin{claim}[The reflexive constraint]
The ethics of reticence forbids its own paternalistic application. Because the wrong is the determination of another's generative domain, any attempt to diagnose her effacement and direct her restoration is a further determination and so a further effacement---the smothering that care ethics' own critics located within devotion. The obligation is therefore strictly to vacate one's overreach and never to manage the other's recovery; and where the question is whether she has been effaced, the question must be returned to her, on the second-personal authority that the claim was always hers to make. An ethic of non-determination that determined, on the other's behalf, that she had been determined would refute itself in the performance.
\end{claim}
